%%%%%%%%%%%%%%%%%%%%%%%%%%%%%%%%%%%%%%%%%%%%%%%%%%%%%%%%%%%%%%%%%%%%%%%
% Thema: AIN Poster-Hintergrund "Labor für mobile Robotik"
% Autor: Jürgen Keppler
% Datum: 18.05.2017
%%%%%%%%%%%%%%%%%%%%%%%%%%%%%%%%%%%%%%%%%%%%%%%%%%%%%%%%%%%%%%%%%%%%%%%
%
\documentclass[a1,landscape]{a0poster}

\usepackage[utf8]{inputenc}
\usepackage[german]{babel}
\usepackage{arev}
\usepackage[T1]{fontenc}
\usepackage{wrapfig}
\usepackage[framemethod=TikZ]{mdframed}
\usepackage{lipsum}
\usepackage{fontspec}
\usepackage{multicol}
\usepackage{adjustbox}
\usepackage{amsmath}
\usepackage{amsfonts}
\usepackage{amssymb}
\usepackage{pgffor}
\usepackage[absolute,overlay]{textpos}
\usepackage{graphicx}
\usepackage{wrapfig}
\usepackage{tikz}
\usetikzlibrary{angles,quotes,babel,plotmarks}
\usetikzlibrary{plotmarks}


\author{Lorenz Bung}
\title{Praxissemester Poster}

\setmainfont[Path=fonts/, BoldFont=swis721-bold.ttf]{swis721.ttf}

\xdefinecolor{color_Black}{rgb}{0.0,0.0,0.0}
\xdefinecolor{color_Raster}{rgb}{1.0,1.0,1.0}
\xdefinecolor{color_White}{rgb}{1.0,1.0,1.0}
\xdefinecolor{color_Teal}{cmyk}{1.0,0.0,0.50,0.0}
\xdefinecolor{color_SoftBlue}{cmyk}{0.10,0.0,0.0,0.10}
\xdefinecolor{color_DarkBlue}{cmyk}{0.80,0.62,0.42,0.45}

\pagecolor{color_DarkBlue}

\setlength{\columnsep}{40pt} % Multicols Spaltenabstand

% No Paragraph Indentation
\setlength{\parindent}{0pt}
\setlength{\parskip}{1em}

% Rasterabstand ab DIN A5 bis DIN A1
\def\unit{145.6pt}

% Rasterpunkt Durchmesser ab DIN A5 d = 0,5 mm, ab DIN A2 d = 1,0 mm
%             bei DIN A1 kann auch d = 2,0 mm gewählt werden
\def\rpoint{0.10cm}

\newenvironment{block}{\begin{mdframed}[roundcorner=25pt,innerleftmargin=40pt,innertopmargin=40pt,innerrightmargin=40pt,innerbottommargin=40pt,backgroundcolor=color_SoftBlue,linecolor=color_SoftBlue]}{\end{mdframed}}

% Width: 2384 Height: 1684
\begin{document}

\begin{textblock*}{2184pt}(100pt,114pt)
	\begin{tikzpicture}
	\foreach \x in {0,1,2,3,4,5,6,7,8,9,10,11,12,13,14,15}
	\foreach \y in {0,1,2,3,4,5,6,7,8,9,10}
	{
		\fill[color_Raster] (\x*\unit,\y*\unit) circle(0.5*\rpoint);
	}
	
	% Dekoration im Kopfbereich
	\draw[color=color_Teal, line width=\rpoint*2.0, cap=round]
	(\unit*2,\unit*10) -- (\unit*14,\unit*9);
	
	\draw[color=color_White, line width=\rpoint*3.0, cap=round]
	(\unit*2,\unit*9) -- (\unit*0,\unit*8);

	\fill[color=color_SoftBlue] (\unit*4,\unit*9) circle[radius=\rpoint*8.0]; 
	
	\draw[color=color_Teal, line width=\rpoint*2.0, cap=round]
	(\unit*2,\unit*8) -- (\unit*11,\unit*8);
	
	\draw[color=color_White, line width=\rpoint*1.0, cap=round]
	(\unit*5,\unit*9) -- (\unit*7,\unit*10);
	
	\draw[color=color_SoftBlue, line width=\rpoint*3.0, cap=round]
	(\unit*15,\unit*9) -- (\unit*13,\unit*8);

	\draw[color=color_Teal, line width=\rpoint*15.0, cap=round]
	(\unit*1,\unit*5) -- (\unit*2,\unit*7);
	
	\fill[color=color_Teal] (\unit*4,\unit*2)   circle[radius=100.0pt];
	\fill[color=color_White] (\unit*14,\unit*1)   circle[radius=75.0pt];
	\fill[color=color_SoftBlue] (\unit*0.5,\unit*0.5)   circle[radius=50.0pt];
	
	\end{tikzpicture}
\end{textblock*}

\begin{textblock*}{2100pt}(142pt,450pt)	
	\begin{multicols}{2}
		\begin{block}
			\begin{wrapfigure}{r}{0.4\textwidth}
				\includegraphics[width=\linewidth]{media/isys.png}
			\end{wrapfigure}
			\textbf{Das Unternehmen}\\
			Die isys vision GmbH und Co. KG ist ein mittelständisches Unternehmen mit ca. 20 Mitarbeitern.
			An ihrem Standort in Freiburg im Breisgau ist isys vision im Bereich der industriellen Bildverarbeitung sowie Automatisierung tätig.\\
			Die Produktpalette umfasst sowohl Soft- und Hardware, welche auf die indivituellen Wünsche der Kunden zugeschnitten ist,
			als auch durch grafische Programmierung konfigurierbare Lösungen.
		\end{block}
		\vfill
		\begin{block}
			\textbf{Einrichtung Intranetserver}\\
			Neben der Webentwicklung war ich einem breiten Aufgabenfeld ausgesetzt. Unter anderem richtete ich auch einen Intranetserver mit RAID5 ein.\\
			Weiterhin erstellte ich Datenbanken, um anschließend Jira und Confluence (Tools für das Production Management) zu installieren.
		\end{block}
		\vfill
		\begin{block}
			\begin{wrapfigure}{r}{0.4\textwidth}
				\includegraphics[width=\linewidth]{media/mikado-roboter.jpg}
			\end{wrapfigure}
			\textbf{Robotikprojekt Mikado}\\
			Mikado ist ein Robotikprojekt, das auf das sogenannte ``Bin-Picking'' ausgelegt ist. Dabei geht es darum, unsortierte Teile zu erkennen und korrekt zu greifen.
			Zur Erkennung der Teile wird eine Punktwolke aufgenommen, in welcher dann die Teile gesucht werden.\\
			Die Software liefert dem Endkunden eine grafische Oberfläche, in welcher die Funktionsweise des Roboters mithilfe von grafischer Programmierung individuell geändert weden kann.
		\end{block}
		\vfill
		\begin{block}
			\begin{wrapfigure}{r}{0.4\textwidth}
				\includegraphics[width=\linewidth]{media/webpreview.png}
			\end{wrapfigure}
			\textbf{WebPreview}\\
			Ein Kernbereich meiner Tätigkeit umfasste die Erstellung einer webbasierten Simulation, welche zur Demonstration der gefundenen Teile genutzt werden sollte. Die Vorschau ist mithilfe von JSON-Daten vollständig konfigurierbar und es können Punktwolken und Meshes zur Darstellung angegeben werden.\\
			Neben Technologien wie Javascript, CSS und HTML wurden die Bibliotheken jQuery und THREE.js zur Animation genutzt.
		\end{block}
		\vfill
		\begin{block}
			\begin{wrapfigure}{r}{0.2\textwidth}
				\includegraphics[width=\linewidth]{media/webremote-info.png}
			\end{wrapfigure}
			\textbf{WebRemote}\\
			Eine weitere Aufgabe war die Erstellung einer mobilen Webseite, welche über ein Smartphone zur Steuerung des Roboters verwendet werden sollte.\\
			Neben der Entwicklung eines passenden responsive Designs stellte auch die Verbindung mit dem Roboter eine Herausforderung dar.\\\\
			ROSlibJS ist eine Bibliothek, welche zur Kommunikation mit dem Roboter verwendet wird. So konnte eine Steuerung zum Starten, Pausieren und schrittweisen Ablauf des Programmes erstellt werden.
		\end{block}
	\end{multicols}
\end{textblock*}

\textblockcolour{color_DarkBlue}

\begin{textblock*}{2100pt}(142pt,1605pt)
	\color{white}{Lorenz Bung \hfill lorenz.bung@htwg-konstanz.de \hfill AIN 4}
\end{textblock*}

\begin{textblock*}{1300pt}(415pt,320pt)
	\Huge{\color{white}{Praxissemester bei isys vision GmbH \& Co. KG}}
\end{textblock*}

\begin{textblock*}{375pt}(1950pt,114pt)
	\includegraphics[scale=4.0]{media/HTWGSchriftzug.eps}
\end{textblock*}

\begin{textblock*}{200pt}(88pt,90pt)
	\includegraphics [scale=2.44]{media/HTWGLogo.eps}
\end{textblock*}

\begin{textblock*}{500pt}(1800pt,1520pt)
	\huge{\color{white}{Angewandte Informatik}}
\end{textblock*}

\begin{textblock*}{200pt}(1800pt,1495pt)
	\color{white}{Studiengang}
\end{textblock*}

\end{document}